% ============================================================
% ch03_database_design.tex
% PostgreSQL Schema + MongoDB WAL Design
% ============================================================

\chapter{PostgreSQL Database Design}

\section{Schema Overview}

The PostgreSQL database \texttt{bank\_recovery\_db} contains six tables.
The \texttt{old\_balance} and \texttt{new\_balance} columns in
\texttt{Transaction\_Record} are critical --- they mirror the WAL log fields
and enable recovery directly from the relational database when needed.

\section{Table Definitions}

\subsection{Branch Table}

\begin{lstlisting}[style=sqlstyle, caption={Branch Table}]
CREATE TABLE Branch (
    branch_id      SERIAL       PRIMARY KEY,
    branch_name    VARCHAR(100) NOT NULL,
    branch_city    VARCHAR(50)  NOT NULL,
    branch_address TEXT,
    contact_number VARCHAR(20),
    created_at     TIMESTAMP    DEFAULT CURRENT_TIMESTAMP,
    CONSTRAINT chk_branch_name CHECK (TRIM(branch_name) <> '')
);
\end{lstlisting}

\subsection{Customer Table}

\begin{lstlisting}[style=sqlstyle, caption={Customer Table}]
CREATE TABLE Customer (
    customer_id    SERIAL       PRIMARY KEY,
    first_name     VARCHAR(50)  NOT NULL,
    last_name      VARCHAR(50)  NOT NULL,
    email          VARCHAR(100) UNIQUE NOT NULL,
    contact_number VARCHAR(20),
    address        TEXT,
    date_of_birth  DATE,
    created_at     TIMESTAMP    DEFAULT CURRENT_TIMESTAMP,
    CONSTRAINT chk_email CHECK (email LIKE '%@%.%'),
    CONSTRAINT chk_dob   CHECK (date_of_birth < CURRENT_DATE)
);
\end{lstlisting}

\subsection{Account Table}

\begin{lstlisting}[style=sqlstyle, caption={Account Table --- Core Banking Entity}]
CREATE TABLE Account (
    account_id     SERIAL        PRIMARY KEY,
    customer_id    INT           NOT NULL,
    branch_id      INT           NOT NULL,
    account_number VARCHAR(20)   UNIQUE NOT NULL,
    account_type   VARCHAR(30)   DEFAULT 'SAVINGS',
    balance        DECIMAL(12,2) DEFAULT 0.00,
    status         VARCHAR(20)   DEFAULT 'ACTIVE',
    created_at     TIMESTAMP     DEFAULT CURRENT_TIMESTAMP,
    CONSTRAINT fk_account_customer
        FOREIGN KEY (customer_id) REFERENCES Customer(customer_id)
        ON DELETE RESTRICT,
    CONSTRAINT fk_account_branch
        FOREIGN KEY (branch_id) REFERENCES Branch(branch_id)
        ON DELETE RESTRICT,
    CONSTRAINT chk_account_type
        CHECK (account_type IN ('SAVINGS','CURRENT','FIXED_DEPOSIT')),
    CONSTRAINT chk_account_status
        CHECK (status IN ('ACTIVE','INACTIVE','FROZEN','CLOSED')),
    CONSTRAINT chk_balance CHECK (balance >= 0.00)
);
\end{lstlisting}

\subsection{Transaction\_Record Table}

\begin{lstlisting}[style=sqlstyle, caption={Transaction\_Record --- WAL-mirroring balance columns}]
CREATE TABLE Transaction_Record (
    transaction_id     SERIAL        PRIMARY KEY,
    account_id         INT           NOT NULL,
    transaction_type   VARCHAR(20),
    amount             DECIMAL(12,2) NOT NULL CHECK (amount > 0),
    old_balance        DECIMAL(12,2) NOT NULL, -- Used in UNDO
    new_balance        DECIMAL(12,2) NOT NULL, -- Used in REDO
    transaction_status VARCHAR(20)   DEFAULT 'PENDING',
    transaction_time   TIMESTAMP     DEFAULT CURRENT_TIMESTAMP,
    notes              TEXT,
    CONSTRAINT fk_transaction_account
        FOREIGN KEY (account_id) REFERENCES Account(account_id),
    CONSTRAINT chk_txn_type
        CHECK (transaction_type IN
              ('DEPOSIT','WITHDRAWAL','TRANSFER_IN','TRANSFER_OUT')),
    CONSTRAINT chk_txn_status
        CHECK (transaction_status IN
              ('PENDING','COMMITTED','ROLLED_BACK'))
);
\end{lstlisting}

\begin{conceptbox}
\textbf{Key Design Decision:} \texttt{old\_balance} and \texttt{new\_balance}
are stored in \texttt{Transaction\_Record} and also in the MongoDB WAL log.
This redundancy is intentional --- it means recovery can use either source
to identify what to undo or redo.
\end{conceptbox}

\subsection{Employee Table}

\begin{lstlisting}[style=sqlstyle, caption={Employee Table}]
CREATE TABLE Employee (
    employee_id    SERIAL      PRIMARY KEY,
    branch_id      INT         NOT NULL,
    first_name     VARCHAR(50) NOT NULL,
    last_name      VARCHAR(50) NOT NULL,
    role           VARCHAR(50) DEFAULT 'TELLER',
    email          VARCHAR(100) UNIQUE,
    hire_date      DATE,
    CONSTRAINT fk_employee_branch
        FOREIGN KEY (branch_id) REFERENCES Branch(branch_id),
    CONSTRAINT chk_role
        CHECK (role IN ('MANAGER','TELLER','LOAN_OFFICER','CASHIER'))
);
\end{lstlisting}

\subsection{Loan Table}

\begin{lstlisting}[style=sqlstyle, caption={Loan Table}]
CREATE TABLE Loan (
    loan_id       SERIAL        PRIMARY KEY,
    customer_id   INT           NOT NULL,
    account_id    INT           NOT NULL,
    loan_type     VARCHAR(50)   DEFAULT 'PERSONAL',
    loan_amount   DECIMAL(12,2) NOT NULL CHECK (loan_amount > 0),
    interest_rate DECIMAL(5,2)  NOT NULL,
    start_date    DATE,
    end_date      DATE,
    loan_status   VARCHAR(20)   DEFAULT 'ACTIVE',
    CONSTRAINT fk_loan_customer
        FOREIGN KEY (customer_id) REFERENCES Customer(customer_id),
    CONSTRAINT fk_loan_account
        FOREIGN KEY (account_id) REFERENCES Account(account_id),
    CONSTRAINT chk_loan_type
        CHECK (loan_type IN ('HOME','CAR','PERSONAL','EDUCATION','BUSINESS')),
    CONSTRAINT chk_loan_status
        CHECK (loan_status IN ('ACTIVE','CLOSED','DEFAULTED'))
);
\end{lstlisting}

\section{Sample Data}

\begin{lstlisting}[style=sqlstyle, caption={Sample INSERT statements}]
-- Branch
INSERT INTO Branch (branch_name, branch_city, branch_address)
VALUES ('Main Branch', 'Lahore', '12 Mall Road, Lahore');

-- Customer
INSERT INTO Customer (first_name, last_name, email, date_of_birth)
VALUES ('Ali', 'Raza', 'ali.raza@gmail.com', '1990-03-15');

-- Account
INSERT INTO Account (customer_id, branch_id, account_number,
                     account_type, balance)
VALUES (1, 1, 'ACC-00001', 'SAVINGS', 10000.00);

-- Transaction with both balance fields (for recovery)
INSERT INTO Transaction_Record
    (account_id, transaction_type, amount,
     old_balance, new_balance, transaction_status, notes)
VALUES (1, 'DEPOSIT', 2000.00, 8000.00, 10000.00,
        'COMMITTED', 'Initial deposit');
\end{lstlisting}

% ────────────────────────────────────────────────────────────
\chapter{MongoDB WAL Design}

\section{Database and Collections}

MongoDB database: \texttt{bank\_recovery\_logs}

\begin{table}[H]
\centering
\caption{MongoDB Collections}
\begin{tabularx}{\textwidth}{llX}
\toprule
\textbf{Collection} & \textbf{Written By} & \textbf{Purpose} \\
\midrule
wal\_logs       & MongoLogger (every transaction) & WAL entries — one per operation \\
recovery\_logs  & RecoveryManager                 & Audit trail of UNDO/REDO actions \\
crash\_log      & CrashSimulator                  & Crash event records \\
system\_config  & MongoLogger                     & log\_id auto-increment counter \\
\bottomrule
\end{tabularx}
\end{table}

\section{wal\_logs Document Structure}

\begin{lstlisting}[style=jsonstyle, caption={WAL Log Document --- UNCOMMITTED state (after crash)}]
{
  "log_id":      8,
  "session_id":  "crash_a1b2c3",
  "account_id":  1,
  "txn_type":    "WITHDRAWAL",
  "old_balance": 10000.00,
  "new_balance":  7000.00,
  "amount":       3000.00,
  "status":      "UNCOMMITTED",
  "timestamp":   "2026-05-20T08:00:00Z",
  "notes":       "CRASH_DEMO: will not commit"
}
\end{lstlisting}

\begin{lstlisting}[style=jsonstyle, caption={WAL Log Document --- COMMITTED state (normal flow)}]
{
  "log_id":      1,
  "session_id":  "sess_f3a7b1",
  "account_id":  1,
  "txn_type":    "DEPOSIT",
  "old_balance": 8000.00,
  "new_balance": 10000.00,
  "amount":      2000.00,
  "status":      "COMMITTED",
  "timestamp":   "2026-05-20T08:00:00Z",
  "notes":       "Deposit to ACC-00001"
}
\end{lstlisting}

\section{Transaction Status Lifecycle}

\begin{figure}[H]
\centering
\begin{tikzpicture}[
  state/.style = {draw, rounded corners=6pt,
                  minimum width=3cm, minimum height=0.9cm,
                  align=center, font=\small\bfseries},
  arr/.style   = {-{Latex[length=3mm]}, thick},
]
\node[state, fill=yellow!30, draw=orange!70]
  (un) at (0,0)   {UNCOMMITTED};
\node[state, fill=green!20, draw=green!60]
  (co) at (-3,-2.5) {COMMITTED};
\node[state, fill=red!15,  draw=red!60]
  (rb) at (3, -2.5) {ROLLED\_BACK};

\draw[arr] (un) -- node[left, font=\footnotesize]
  {PG commit\\succeeds} (co);
\draw[arr] (un) -- node[right, font=\footnotesize]
  {ROLLBACK or\\crash recovery} (rb);
\draw[arr, dashed, gray] (-0.5, 0.8) -- node[above, font=\footnotesize]
  {WAL log written} (un.north);
\end{tikzpicture}
\caption{WAL Log Status Lifecycle}
\label{fig:status}
\end{figure}

\section{recovery\_logs Document}

\begin{lstlisting}[style=jsonstyle, caption={Recovery Log --- written by RecoveryManager after UNDO}]
{
  "action":      "UNDO",
  "account_id":  1,
  "before_bal":  7000.00,
  "after_bal":   10000.00,
  "reason":      "CRASH_DETECTED",
  "timestamp":   "2026-05-20T08:01:00Z"
}
\end{lstlisting}

\section{Useful MongoDB Queries}

\begin{lstlisting}[style=jsonstyle, caption={Key mongosh queries for demo and viva}]
// Find all UNCOMMITTED logs (crash victims)
db.wal_logs.find({ status: "UNCOMMITTED" })

// Find all COMMITTED logs
db.wal_logs.find({ status: "COMMITTED" })

// Count by status (shows system state at a glance)
db.wal_logs.aggregate([
    { $group: { _id: "$status", count: { $sum: 1 } } }
])

// Show UNDO plan: what to restore for each crash victim
db.wal_logs.find(
    { status: "UNCOMMITTED" },
    { log_id:1, account_id:1, old_balance:1, _id:0 }
)

// View all recovery actions taken
db.recovery_logs.find({}, { _id: 0 })
\end{lstlisting}

\section{Why MongoDB for WAL Logs?}

\begin{table}[H]
\centering
\caption{PostgreSQL vs MongoDB for Log Storage}
\begin{tabularx}{\textwidth}{lXX}
\toprule
\textbf{Criterion} & \textbf{PostgreSQL} & \textbf{MongoDB} \\
\midrule
Schema flexibility & Fixed columns only & Any fields per document \\
Insert speed & Transaction overhead & Lightweight insert \\
Log independence & Affected by PG crash & Independent of PG failures \\
Readability & Binary/SQL format & Human-readable JSON \\
Query language & SQL & Rich document queries \\
Status updates & UPDATE statement & \$set operator \\
\bottomrule
\end{tabularx}
\end{table}
